% !PRINTED 616
\weberpage{616}

Es bleibt nur noch \"ubrig, zu zeigen, wie man $\varphi_2$ der hier gestellten Forderung gem\"ass bestimmen kann. Da zwei Formen, die sich nur durch einen Zahlenfactor von einander unterscheiden, \"aquivalent sind, so ist $\varphi_2$ \"aquivalent mit
\[
c_2\, (t_1 + \omega_2 t_2),
\]
[\S.\,166, (10)] und nehmen wir diese Form f\"ur $\varphi_2$, so ist $b_{1,1} = c_2$. Die Zahl $\omega_2$ kann man aber unter den mit ihr \"aquivalenten immer so w\"ahlen, dass $c_2$ zu $a_{1,1}$ relativ prim wird, und dann ist auch $a_{1,1}$ zu $\varphi_2$ relativ prim.

Ersetzt man n\"amlich $\omega_2$ durch
\[
\frac{p\, \omega_2 + q}{r\, \omega_2 + s}, \quad p s - q r = \pm 1,
\]
so geht $c_2$ nach Bd.\ I, \S.\,122, (13) in
\[
c_2' = -a_2 r^2 + b_2 r s + c_2 s^2
\]
\"uber, und nun kann man \"uber die ganzen rationalen Zahlen $r, s$ so verf\"ugen, dass sie unter einander theilerfremd sind und dass $c_2'$ durch irgend welche gegebene Primzahlen nicht theilbar ist.

Denn ist $n$ eine solche Primzahl, so nehme man, wenn $n$ in $a_2$ und $c_2$, also nicht in $b_2$ aufgeht, $r$ und $s$ durch $n$ untheilbar, wenn $n$ in $a_2$ nicht aufgeht, $s$ durch $n$ theilbar und $r$ durch $n$ untheilbar, und wenn $n$ in $c_2$ nicht aufgeht, $r$ durch $n$ theilbar, $s$ durch $n$ untheilbar. (Sind $a_2$ und $c_2$ beide nicht durch $n$ theilbar, so hat man zwischen den beiden letzteren Annahmen die Wahl.) Diese Forderungen sind f\"ur verschiedene beliebig gegebene Primzahlen $n$ mit einander vertr\"aglich, und dann l\"asst sich $p, q$ aus der Bedingung $p s - q r = \pm 1$ bestimmen.

Diese Gesetze der Composition der quadratischen Irrationalzahlen haben wir hier nur unter der Voraussetzung abgeleitet, dass die Discriminante zugleich die Grundzahl eines quadratischen K\"orpers ist, d.\,h. dass $\Delta$ Stammdiscriminante ist. Dieselben Gesetze gelten aber auch, mit geringen Modificationen, allgemein $^{1)}$.

\noindent\rule{4cm}{0.4pt}

$^{1)}$ \gesperrt{Gauss}, Disquisitiones arithm., Art.\,234 ff. \gesperrt{Dirichlet-Dedekind}, Zahlentheorie. 4. Auflage, \S.\,187.

\begin{center}
\rule{0.2\linewidth}{0.4pt}
\end{center}
